\documentclass[10pt,a4paper]{amsart}
\usepackage[margin=19mm]{geometry}
\usepackage{amsmath,amssymb,mathtools}
\usepackage[scr=rsfs,cal=euler]{mathalfa}
\usepackage{xfrac}
\usepackage[unicode,hidelinks,bookmarks=false]{hyperref}
\newcommand{\ie}{i.e.\ }
\newcommand{\cf}{cf.\ }
\newcommand{\resp}{respectively\ }
\newcommand{\Subsection}{\subsection}
\providecommand{\nobreakdash}{\nobreak\hbox{-}}
\setlength{\emergencystretch}{2em}
\multlinegap0pt
\usepackage{bm}
\usepackage{amssymb}
\usepackage{mathtools}
\usepackage[shortlabels]{enumitem}
\usepackage{xcolor}
\newtheorem{proposition}{Proposition}
\newcommand{\e}{\mathrm{e}}
\title{Rational arrival processes with strictly positive densities need not be Markovian}
\author{Oscar Peralta}
\date{Source: arXiv:2603.28047v1 (CC BY 4.0)}
\begin{document}
\maketitle

\noindent\emph{Source and licence.} Rational arrival processes with strictly positive densities need not be Markovian. Version arXiv:2603.28047v1 (CC BY 4.0), distributed by arXiv under the Creative Commons Attribution 4.0 International licence, http://creativecommons.org/licenses/by/4.0/, as recorded in the arXiv licence metadata for this exact version and shown on the abstract page https://arxiv.org/abs/2603.28047v1. Full licence notice: \texttt{licenses/arxiv-2603.28047-CC-BY-4.0.txt}.\par\smallskip

\noindent\textit{Editorial scope.} Counted unit: the article's proposition prop:positivity (source0.tex lines 340--350). The article's complete original proof of it is delivered with it (source0.tex lines 351--452, one \texttt{proof} environment of 102 lines). No mathematics is written, repaired or re-derived: the statement and the proof are the article's own lines, in the article's own notation, and no retained line is substituted or re-flowed.  Retained uncounted as the source-local context the proof needs: lines 184--199.  Nothing was omitted from any retained span, so the package carries no omission record.  No retained line contains a citation, so the wrapper carries no bibliography and no bibliography entry.  No retained line contains a figure environment, an image-inclusion command or drawing code, so no image is delivered and the package declares no asset dependency.  Editorial formatting only, in addition: an AMS \texttt{amsart} preamble that restates the article's own declarations and macros used by the retained lines, namely the environments \texttt{proposition} on separate counters, together with the standard packages those lines need.  The retained environments print in this document as Proposition 1 in order of appearance, whereas the article prints the same items with the numbering of its own full document.  The complete native source of the article as distributed in the arXiv e-print for this exact version is bundled unchanged as \texttt{sources/197404/source0.tex}.
\begin{equation}\label{eq:density-formula}
f_k(t_1,\dots,t_k)
=
\e^{-T}
\left[
1-\varepsilon
+
\varepsilon\,(1,0)
\left(
\sum_{j=1}^{k-1}
\e^{-T_j}(2\bm{R}_\varphi^{j-1}-\bm{R}_\varphi^j)\bm{1}_2
+
2\e^{-T}\bm{R}_\varphi^{k-1}\bm{1}_2
\right)
\right],
\end{equation}

\begin{proposition}\label{prop:positivity}
There is a constant $M=M(\varphi)$ such that, whenever
\[
0<\varepsilon<\frac{1}{M+1},
\]
we have
\[
f_k(t_1,\dots,t_k)>0
\qquad\text{for every }k\ge 1\text{ and every }t_1,\dots,t_k\ge 0.
\]
\end{proposition}

\begin{proof}
Fix $k_0\ge 1$. By \eqref{eq:density-formula},
\[
f_{k_0}(t_1,\dots,t_{k_0})
=
\e^{-T}
\left[
1-\varepsilon
+
\varepsilon
\left(
\sum_{j=1}^{{k_0}-1}
\e^{-T_j}(1,0)(2\bm{R}_\varphi^{j-1}-\bm{R}_\varphi^j)\bm{1}_2
+
2\e^{-T}(1,0)\bm{R}_\varphi^{{k_0}-1}\bm{1}_2
\right)
\right].
\]
Set
\[
w_j:=(1,0)(2\bm{R}_\varphi^{j-1}-\bm{R}_\varphi^j)\bm{1}_2
\quad (1\le j\le {k_0}-1),
\qquad
w_{k_0}:=2(1,0)\bm{R}_\varphi^{{k_0}-1}\bm{1}_2,
\]
and let
\[
W_m:=\sum_{j=1}^m w_j.
\]
A direct telescoping computation gives
\[
W_m
\;=\; (1,0)\left( 2\sum_{j=0}^{m-1}\bm{R}_\varphi^j - \sum_{j=1}^{m}\bm{R}_\varphi^j \right)\bm{1}_2
=
(1,0)\left(\bm{I}_2-\bm{R}_\varphi^m+\sum_{\ell=0}^{m-1}\bm{R}_\varphi^\ell\right)\bm{1}_2,
\qquad 1\le m\le {k_0}-1,
\]
and
\[
W_{k_0}
\;=\; W_{k_0-1} + 2(1,0)\bm{R}_\varphi^{k_0-1}\bm{1}_2
=
(1,0)\left(\bm{I}_2+\sum_{\ell=0}^{{k_0}-1}\bm{R}_\varphi^\ell\right)\bm{1}_2.
\]

Identify $\mathbb{R}^2$ with $\mathbb{C}$ via $(x,y)\leftrightarrow x+iy$, so that $\bm{R}_\varphi$ acts as multiplication by $e^{i\varphi}$. Any real polynomial $\bm{A}$ in $\bm{R}_\varphi$ corresponds to a complex scalar $z_{\bm{A}}$, and $(1,0)\bm{A}\bm{1}_2 = \operatorname{Re}(z_{\bm{A}}(1+i))$, so $|(1,0)\bm{A}\bm{1}_2|\le\sqrt{2}|z_{\bm{A}}|$. Specifically, for
\[
\bm{A}_m := \bm{I}_2-\bm{R}_\varphi^m+\sum_{\ell=0}^{m-1}\bm{R}_\varphi^\ell,
\quad\mbox{we have}\quad
z_{\bm{A}_m} = 1-e^{im\varphi}+\sum_{\ell=0}^{m-1}e^{i\ell\varphi} = 1-e^{im\varphi}+\frac{1-e^{im\varphi}}{1-e^{i\varphi}}.
\]
Since \(\varphi\in(0,\pi)\), we have \(e^{i\varphi}\neq 1\), so the denominator \(|1-e^{i\varphi}|=2\sin(\varphi/2)\) is strictly positive (dropping the absolute value since \(\sin(\varphi/2)>0\)). By the triangle inequality,
\[
|z_{\bm{A}_m}|
\le
|1-e^{im\varphi}|+\frac{|1-e^{im\varphi}|}{|1-e^{i\varphi}|}
\le
2+\frac{2}{2\sin(\varphi/2)}
=
2+\frac{1}{\sin(\varphi/2)}.
\]
Therefore,
\[
|W_m| \le \sqrt{2}|z_{\bm{A}_m}| \le \sqrt{2}\left(2+\frac{1}{\sin(\varphi/2)}\right).
\]
The same logic applies to \(W_{k_0}\). Letting \(\bm{A}_{k_0} := \bm{I}_2+\sum_{\ell=0}^{k_0-1}\bm{R}_\varphi^\ell\), its complex counterpart \(z_{\bm{A}_{k_0}}\) gives
\[
|W_{k_0}|
\le
\sqrt{2}|z_{\bm{A}_{k_0}}|
=
\sqrt{2}\left|1+\frac{1-e^{ik_0\varphi}}{1-e^{i\varphi}}\right|
\le
\sqrt{2}\left(1+\frac{1}{\sin(\varphi/2)}\right).
\]
Setting \(M:=\sqrt{2}\left(2+\frac{1}{\sin(\varphi/2)}\right)\), hence
\[
|W_m|\le M
\qquad\text{for all }1\le m\le k_0.
\]

Now write
\[
\sum_{j=1}^{k_0} e^{-T_j}w_j
=
\sum_{j=1}^{{k_0}-1}W_j\bigl(e^{-T_j}-e^{-T_{j+1}}\bigr)+W_{k_0}e^{-T_{k_0}}.
\]
Since \(0\le e^{-T_{k_0}}\le\cdots\le e^{-T_1}\le1\), it follows that
\[
\left|\sum_{j=1}^{k_0} e^{-T_j}w_j\right|
\le
M\sum_{j=1}^{{k_0}-1}(e^{-T_j}-e^{-T_{j+1}})+Me^{-T_{k_0}}
=
Me^{-T_1}
\le M.
\]
Therefore
\[
f_{k_0}(t_1,\dots,t_{k_0})\ge e^{-T}\bigl(1-\varepsilon-\varepsilon M\bigr),
\]
which is strictly positive as soon as \(\varepsilon<1/(M+1)\).
\end{proof}

\end{document}
